%! Unit Opener: Functions as Models — Unit 2, Lesson 2.0 — Exit Ticket
\documentclass[10pt]{article}
\usepackage{atda-article}
\usepackage{atda-boxes}
\newcommand{\TallMath}[1]{$\displaystyle #1\rule[-1.4em]{0pt}{3.2em}$}
\setlength{\workrowsep}{8pt}

\begin{document}

\pageheader{Unit 2, Lesson 2.0}{Exit Ticket}
% NAMESTRIP: no \namedateperiod here. The student writes their name once, on the
% cover this component is stapled behind. See references/conventions.md.

\vspace{0.15in}

\begin{notesbox}{Reading a Graph --- On Your Own}
\small
Notes closed. Three items. A soccer ball is kicked straight up from the ground; $h$ is its height
in feet, $t$ seconds after the kick.

\begin{center}
\begin{tikzpicture}
\begin{axis}[
    width=9.8cm, height=4.6cm,
    xlabel={$t$ (seconds)}, ylabel={$h$ (feet)},
    xmin=0, xmax=4.4, ymin=0, ymax=75,
    xtick={0,1,2,3,4}, ytick={0,16,32,48,64},
    grid=both, grid style={linegray!45}, axis lines=left,
    tick label style={font=\scriptsize}, label style={font=\scriptsize}]
  \addplot[royal, very thick, domain=0:4, samples=60]{-16*x^2+64*x};
\end{axis}
\end{tikzpicture}
\end{center}

\begin{enumerate}[leftmargin=1.6em, itemsep=8pt, topsep=4pt]

\item State the domain and the range of this function in this situation, with units.

\writelines{1}

\item Give the absolute maximum: its value, where it happens, and what it means for the ball.

\writelines{1}

\item A classmate says, ``The domain of every function is all real numbers --- you can always put
any number in.'' Do you agree? Answer in one or two sentences, using this graph.

\writelines{2}

\end{enumerate}
\end{notesbox}

\end{document}
